\chapter{Pemecahan Masalah}

\section{Masalah Instalasi}

\subsection{Perintah \texttt{pakai} tidak ditemukan}

\textbf{Gejala:} \texttt{bash: pakai: command not found}

\textbf{Penyebab:} Direktori instalasi tidak ada di \texttt{\$PATH}.

\textbf{Solusi:}
\begin{lstlisting}[style=terminal]
# Cari lokasi binary
$ which pakai || find ~/go ~/go/bin ~/.local/bin -name pakai 2>/dev/null

# Tambahkan ke PATH (sesuaikan path yang ditemukan)
$ echo 'export PATH="$PATH:$HOME/go/bin"' >> ~/.bashrc
$ source ~/.bashrc
\end{lstlisting}

Jika menggunakan \textbf{asdf}:
\begin{lstlisting}[style=terminal]
$ asdf reshim golang
\end{lstlisting}

\subsection{Build gagal: versi Go terlalu lama}

\textbf{Gejala:} \texttt{requires go >= 1.21}

\textbf{Solusi:} Perbarui Go ke versi 1.21 atau lebih baru.
\begin{lstlisting}[style=terminal]
$ go version        # cek versi saat ini
$ asdf install golang latest   # update via asdf
\end{lstlisting}

\section{Masalah Daemon}

\subsection{Daemon gagal dimulai: port sudah digunakan}

\textbf{Gejala:} \texttt{port 7731 is already in use by another process}

\textbf{Solusi:}
\begin{lstlisting}[style=terminal]
# Cek proses yang menggunakan port
$ lsof -i :7731 || ss -tlnp | grep 7731

# Hentikan daemon lama
$ pakai daemon stop

# Atau gunakan port berbeda
$ pakai config set daemon.port 7732
$ pakai daemon start
\end{lstlisting}

\subsection{Daemon berhenti tiba-tiba}

\textbf{Solusi:} Periksa log systemd:
\begin{lstlisting}[style=terminal]
$ journalctl --user -u pakai -n 50
\end{lstlisting}

\subsection{Status daemon menunjukkan data lama (\textit{stale})}

\textbf{Penyebab:} Interval polling terlalu jarang atau koneksi internet bermasalah.

\textbf{Solusi:}
\begin{lstlisting}[style=terminal]
# Perkecil interval polling
$ pakai config set daemon.poll_interval 15
$ pakai daemon stop
$ pakai daemon start
\end{lstlisting}

\section{Masalah Penyedia Claude}

\subsection{Claude stats cache tidak ditemukan}

\textbf{Gejala:} \texttt{Claude stats cache not found: ensure Claude Code is running}

\textbf{Penyebab:} Berkas \texttt{\textasciitilde/.claude/stats-cache.json} tidak ada.

\textbf{Solusi:} Pastikan Claude Code CLI telah diinstal dan pernah dijalankan:
\begin{lstlisting}[style=terminal]
$ ls ~/.claude/stats-cache.json
$ claude --version   # verifikasi instalasi
\end{lstlisting}

\subsection{Kredensial Claude tidak ditemukan}

\textbf{Gejala:} \texttt{Claude credentials not found: ensure Claude Code is installed}

\textbf{Solusi:}
\begin{lstlisting}[style=terminal]
$ claude login
\end{lstlisting}

\subsection{Persentase Claude tidak tersedia}

\textbf{Gejala:} \texttt{live percentage unavailable}

\textbf{Penyebab:} API OAuth gagal, tetapi data lokal tersedia. PakAI tetap
menampilkan data cache.

\textbf{Solusi:} Periksa koneksi internet. Jika masalah berlanjut:
\begin{lstlisting}[style=terminal]
$ claude logout && claude login
\end{lstlisting}

\section{Masalah Penyedia OpenAI Codex}

\subsection{Kredensial Codex tidak ditemukan}

\textbf{Gejala:} \texttt{Pi Codex credentials not found}

\textbf{Solusi:} Pilih sumber Pi atau Codex CLI di Settings, lalu login pada
sumber tersebut.

\subsection{Permintaan penggunaan Codex gagal}

\textbf{Gejala:} \texttt{codex usage request failed: status 401}

\textbf{Penyebab:} Token akses kedaluwarsa.

\textbf{Solusi:} Sambungkan ulang sumber Codex yang dipilih.

\subsection{Kesalahan sertifikat TLS}

\textbf{Gejala:} \texttt{tls: failed to verify certificate: x509: certificate signed by unknown authority}

\textbf{Penyebab:} Proxy perusahaan atau VPN yang melakukan \textit{TLS inspection}.

\textbf{Solusi:} Tambahkan sertifikat CA ke kepercayaan sistem, atau nonaktifkan
sementara penyedia yang bermasalah:
\begin{lstlisting}[style=terminal]
$ pakai config set provider.openai.enabled false
\end{lstlisting}

\section{Masalah Penyedia OpenCode}

\subsection{Database OpenCode tidak ditemukan}

\textbf{Gejala:} \texttt{opencode database not found: ensure opencode is installed}

\textbf{Penyebab:} OpenCode belum pernah digunakan atau database belum terbentuk.

\textbf{Solusi:} Gunakan OpenCode minimal satu kali sehingga database terbentuk:
\begin{lstlisting}[style=terminal]
$ ls ~/.local/share/opencode/
# Harus ada: opencode-stable.db atau opencode.db
\end{lstlisting}

\subsection{Database OpenCode sibuk (\textit{locked})}

\textbf{Gejala:} \texttt{database is busy: another process may be writing}

\textbf{Penyebab:} OpenCode sedang aktif menulis ke database.

\textbf{Solusi:} Tunggu beberapa detik, lalu coba lagi. PakAI akan berhasil
pada polling berikutnya secara otomatis.

\section{Masalah Tampilan}

\subsection{Ikon tidak muncul di tmux}

\textbf{Penyebab:} Font terminal tidak mendukung Nerd Font.

\textbf{Solusi:} Instal dan atur font Nerd Font di terminal dan tmux. Contoh untuk
terminal berbasis X11:
\begin{lstlisting}[style=terminal]
# Pasang font (Ubuntu/Debian)
$ sudo apt install fonts-jetbrains-mono

# Atau unduh dari nerdfonts.com
\end{lstlisting}

\subsection{Waybar tidak memperbarui data}

\textbf{Penyebab:} Interval Waybar atau daemon terlalu panjang.

\textbf{Solusi:} Sesuaikan interval di konfigurasi Waybar:
\begin{lstlisting}[style=toml]
"custom/pakai": {
    "interval": 15
}
\end{lstlisting}

\section{Diagnosis Umum}

Jika masalah tidak tercantum di atas, gunakan langkah diagnosis berikut:

\begin{lstlisting}[style=terminal]
# 1. Cek status daemon
$ pakai daemon status

# 2. Debug penyedia yang bermasalah
$ pakai provider debug claude
$ pakai provider debug opencode
$ pakai provider debug openai

# 3. Cek log daemon
$ journalctl --user -u pakai -n 100

# 4. Cek koneksi ke daemon
$ curl -s http://127.0.0.1:7731/health

# 5. Reset daemon
$ pakai daemon stop && pakai daemon start
\end{lstlisting}
